\documentclass[tikz,border=6pt]{standalone}
\usepackage{peabra}
\begin{document}
\begin{tikzpicture}[x=1cm,y=1cm, every node/.append style={font=\scriptsize},
  el/.style={draw=peKek, line width=1pt}, pn/.style={draw=peZold, line width=1.1pt, dash pattern=on 4pt off 2pt},
  mod/.style={doboz=17mm, minimum height=9mm, inner sep=2pt}]
  \node[sotetdoboz=22mm, minimum height=12mm, font=\scriptsize\bfseries] (ecu) at (4.2,0) {EBS-vezérlő\\(elektronikus kör)};
  \node[mod] (fbm) at (0,0) {lábfékmodul\\(1E + 2P)};
  \node[mod] (fa) at (4.2,2.3) {első tengely\\modulátora};
  \node[mod] (ha) at (4.2,-2.3) {hátsó tengely\\modulátora};
  \node[mod] (tcm) at (8.4,-2.3) {pótkocsi-\\vezérlő modul};
  \node[szurkedoboz=20mm, minimum height=9mm, inner sep=2pt] (veh) at (8.4,0) {jármű-CAN,\\vezetéstámogatás};
  \draw[el] (fbm) -- node[above, pos=0.55, font=\tiny]{elektromos fékigény} (ecu);
  \draw[el] (ecu) -- node[right]{fék-CAN} (fa); \draw[el] (ecu) -- (ha); \draw[el] (ecu.east) -- ++(0.6,0) |- (tcm.west);
  \draw[el] (veh) -- node[above, font=\tiny]{külső fékigény} (ecu);
  \draw[pn] (fbm.north) |- node[pos=0.75, above, text=peZold]{1. pneumatikus kör} (fa.west);
  \draw[pn] (fbm.south) |- node[pos=0.75, below, text=peZold]{2. pneumatikus kör} (ha.west);
  \draw[pn] (ha.east) -- (tcm.west);
  % jelmagyarázat
  \draw[el] (0,-4.0) -- ++(0.8,0) node[right, black]{elektronikus vezérlés (normál üzem)};
  \draw[pn] (4.6,-4.0) -- ++(0.8,0) node[right, black]{pneumatikus vészkör (tartalék)};
  % automatizált vezetés
  \node[kiemeltdoboz=40mm, anchor=north west, align=left, font=\scriptsize] at (10.0,2.4) {\textbf{automatizált vezetéshez:}\\a pneumatikus vészkör nem viszi át a külső fékigényt, ezért egy második, független elektronikus kör (redundáns vezérlő és tápellátás) kell: a fék hibatűrővé válik};
\end{tikzpicture}
\end{document}
